\section{Incidencia excepcional y transporte de formas}
\label{nuclear:4}
\subsection{Generación del registro y conservación de incidencia}

\subsubsection{Registro dodecafásico de la historia enriquecida}

La sección~\ref{subsec:k-registro-integral} define

\[
\mathcal R_{12}(x)=((E_m,\tau_m,\sigma_m,\kappa_m,\Lambda_m))_{m=1}^{12}.
\]

Conserva ventana, subruta, orientación, acarreo y libro de incidencias. La sección~\ref{subsec:k-terminal} fija

\[
x_{\rm term}=(\mathcal U_{108}\cdots\mathcal U_1)x_{\rm seed},
\quad \mathcal U_t=\mathsf{Upd}_t\mathsf{Tra}_t\mathsf{Sel}_t,
\]

\[
R_{\rm sgn}=\Pi_H\mathcal E_{108}^{90,120}\mathcal R_{12}:
\mathcal X_{\rm TPK}^{\rm term,mem}\to\mathbb Z^{12}.
\]

La reconstrucción algebraica siguiente comienza en esa salida tipada. No sustituye la evaluación de los eventos anteriores ni obtiene sus coordenadas imponiendo \(K\) o alfa como objetivo.

\subsubsection{Normalización integral y normalizaciones isométricas}

Para las órbitas de paso tres \((1,4,7,10)\), \((2,5,8,11)\), \((3,6,9,12)\), el lema~\ref{lem:k-hadamard} fija la matriz simétrica \(H=H_{12}\) y prueba \(H^2=4I\).

\[
\mathcal L_H=\{u\in\mathbb Z^{12}:Hu\equiv0\pmod4\},
\qquad \mathscr K(u)=\tfrac14Hu,
\qquad \mathscr K^{-1}(k)=Hk.
\]

Es un isomorfismo integral entre la subred de registros firmados y \(\mathbb Z^{12}\). El factor \(1/4\) realiza la inversión, no una isometría euclídea.

Como consecuencia inmediata de la misma identidad, \(J_H=H/2\) satisface \(J_H^*J_H=I\): éste es el transporte isométrico real. Si \(u=Hk\), se tiene \(u/2=J_Hk\) y \(\|u\|^2=4\|k\|^2\). Confundir \(H/4\) con \(H/2\) cambiaría el balance por un factor cuatro.

En \cref{exc:entrelazador,exc:isometria}, el diseño construido determina

\[
I:\mathbb R^{12}\to\mathbb R^{132},\qquad
(Iv)(H')=\sum_{p\in H'}v_p,\qquad I^*I=36I_{12}+30J_{12}.
\]

Los coeficientes proceden de las incidencias: cada punto pertenece a 66 hexadas y cada pareja a 30. Sobre \(V=\mathbf1^\perp\) desaparece \(J_{12}\), por lo que \(U_W=I/6:V\to E_{\rm inc}\) es isometría. El factor \(1/6\) es el normalizador de esta pantalla, no el inversor de \(K\).

\subsubsection{Refinamiento de prefijos y particiones de eventos}

La sección~\ref{exc:doble-lectura} construye marcos causales \(f_{j+1}=g_j(x)f_j\) y codificaciones \(Q_n\) del prefijo. Prueba

\[
r_{n,m}Q_n=Q_mp_{n,m},\qquad
\gamma_{A^f}(wf^{-1})=\gamma_A(w)(f^{-1}\oplus f^{-1}).
\]

La prueba usa que una prolongación conserva las palabras, transportes y marcos anteriores. El límite inverso produce \(Q_\infty\). La recuperación de la palabra visible se debe a la primera componente de \(\gamma\); el texto distingue expresamente ese alcance de la recuperación de toda memoria profunda y del paso a soportes.

La composición aditiva y afín del registro, descrita en \cref{subsec:k-registro-integral,iv:cont:concatenacion}, prueba una segunda conservación: con contribuciones prospectivas compatibles \((\delta w_t,\delta q_t)\), la inversión del registro es aditiva. Concatenar segmentos o subdividir la misma lista de eventos da el mismo registro final. La versión con transporte no trivial usa el cociclo afín. Conservar el total al cambiar la partición no equivale a afirmar que el registro deja de crecer cuando se añaden nuevos eventos.

\subsection{Transporte K–Witt de normas, operadores y balances}\label{nuc04:kwitt}

Se reúne aquí una formalización de los resultados recuperados, con la prueba algebraica que sigue. Este ensamblaje no atribuye prioridad nueva a sus identidades de partida.

Sean \(P_0=I-J_{12}/12\), \(V=P_0\mathbb R^{12}\), \(W_H=J_HV\) y

\[
B:=U_WJ_H|_{W_H}:W_H\longrightarrow E_{\rm inc}.
\]

Como \(J_H^2=I\), éste es un isomorfismo isométrico entre espacios de dimensión once. Para un registro \(k\), su coordenada firmada normalizada es \(J_Hk\). Si \(\widehat P_0=J_HP_0J_H\),

\[
B\widehat P_0(J_Hk)=U_WP_0k,
\]

\[
\|\widehat P_0(J_Hk)\|^2
=\|U_WP_0k\|^2
=\sum_{p=1}^{12}k_p^2-\frac{(\sum_pk_p)^2}{12}.
\]

No se ha introducido ningún valor de constante para obtener esta igualdad: procede de \(H^2\) y del Gramiano de incidencia. La media eliminada debe conservarse aparte si se quiere recuperar todo \(k\).

Para \(g\in\operatorname{Aut}(\mathcal H)\), sea \(\rho_{12}(g)\) su permutación de puntos y \(\rho_{\rm hex}(g)\) la de hexadas. La acción sobre las coordenadas firmadas es

\[
\widehat\rho_H(g)=J_H\rho_{12}(g)J_H
=\tfrac14H\rho_{12}(g)H.
\]

Es ortogonal, conserva \(W_H\) y también conserva la subred \(\mathcal L_H\): si \(u=Hk\), su imagen es \(H\rho_{12}(g)k\). El entrelazamiento material de \cref{exc:entrelazador} da

\[
B\widehat\rho_H(g)=\rho_{\rm hex}(g)B.
\]

\begin{proof}
sustituir las definiciones, cancelar \(J_H^2\) y aplicar \(U_W\rho_{12}(g)=\rho_{\rm hex}(g)U_W\). Se preservan así el registro integral, el producto interno centrado y la acción incidencial, cada uno en su dominio.
\end{proof}

Para un operador \(A\) ya construido sobre \(W_H\) se obtiene \(A_{\rm inc}=BAB^*\). Por sustitución,

\[
BA=A_{\rm inc}B,\quad
\langle z,Az\rangle=\langle Bz,A_{\rm inc}Bz\rangle.
\]

Por tanto, si un transporte \(C\) sobre \(W_H\) conserva \(A\), \(C^*AC=A\), su transporte \(C_{\rm inc}=BCB^*\) conserva \(A_{\rm inc}\). La condición de que un \(C\) particular sea una publicación de la dinámica TPK permanece en la composición causal que lo produce. La isometría no transforma por sí sola una permutación de incidencia en el ciclo nonádico.

No hace falta imponer conmutación para comparar cartas transportadas. Si \(R\) es unitario, \(C'=RCR^*\) y \(T(C)=(8I+C)/9\), entonces \(T(C')R=RT(C)\) por sustitución. Junto con \(A'=RAR^*\), el balance en una carta se conjuga al de la otra. La conmutación \([R,C]=0\) sólo corresponde al problema adicional de mantener fija esa misma costura \(C\). Esta distinción coincide con la naturalidad reticular de la sección~\ref{nuc04:c3}.

Una acción finita concreta disponible es la estrella seleccionada por la bandera:

\[
\sigma_{B_K}=(1\ 3)(2\ 11)(5\ 7)(8\ 12).
\]

La sección~\ref{exc:estrella} prueba su construcción y la distingue del grupo generado por las 495 estrellas. La estrella individual da \(C_2\); la familia da \(M_{12}\). Transportar la marca y el origen da covariancia entre marcos. Si se pretende una simetría de un marco fijo, el dominio adecuado es su estabilizador, no automáticamente \(M_{12}\) entero. Si además debe conmutar con \(C\), se utiliza el subgrupo que satisfaga la ecuación concreta \([\widehat\rho_H(g),C]=0\).

\subsection{Acción de orden tres sobre las fibras reticulares}\label{nuc04:c3}

La construcción de \cref{nuclear:fibras-app-witt,nuclear:alineamiento-witt} produce doce fibras \(A_2\) etiquetadas por \((i,\mu,\varepsilon)\in\mathbb Z/3\mathbb Z\times\{0,1\}\times\{0,1\}\). El alineamiento \(\lambda(i,\mu,\varepsilon)=4i+2\mu+\varepsilon\pmod{12}\) es biyectivo. Si \(C_{A_2}\) es el Coxeter y \(R\) la reflexión con \(RC_{A_2}R=C_{A_2}^{-1}\), las trivializaciones y marcos transportados fijan

\[
g_b=\iota_b^{-1}P_bC_{A_2}P_b^{-1}\iota_b,
\quad T_{i\mu\varepsilon}=\iota_{\lambda(i,\mu,\varepsilon)}^{-1}P_{\lambda(i,\mu,\varepsilon)}R^\mu,
\]

\[
g_\lambda T_{i\mu\varepsilon}
=T_{i\mu\varepsilon}C_{A_2}^{(-1)^\mu}.
\]

La suma directa \(T_\lambda:W_{\rm APP}\otimes\mathbb C\to\bigoplus_{b=1}^{12}(A_{2,b}\otimes\mathbb C)\) es un isomorfismo \(C_3\)-equivariante. La prueba consiste en la conjugación y en la biyección de las doce etiquetas: conserva par, hoja y órbita. No es una coincidencia de dimensiones.

Las construcciones de \cref{km:palabra-total,km:estabilidad-vecino} utilizan \(q^\star=(1,1,1,1,1,1)\) y su palabra \(c^\star=(q^\star,q^\star A_W)\) de soporte total para orientar los doce bloques. El \(g_c\) resultante tiene orden tres, no tiene vectores fijos no nulos y preserva el pegado de Niemeier. Para el vector incidencial \(v\) de norma cuadrática \(54\),

\[
N_0=\ker(\langle\cdot,v\rangle\bmod3),\qquad
N_v=N_0+\mathbb Z\frac v3,
\]

\[
y_*=(g_cv-v)/3\in N_0,\qquad
z_*=(g_c^{-1}v-v)/3\in N_0.
\]

La prueba identifica sus clases discriminantes como \(c^\star\) y \(-c^\star\), y sus productos con \(v\) como \(-27\). De ahí \(g_cN_0=N_0\) y \(g_c(v/3)=v/3+y_*\), luego \(g_cN_v=N_v\). La estructura conservada incluye el pegado, el núcleo de congruencia y la clase que genera el vecino de Leech.

La prolongación a las 24 orientaciones de soporte total y la elevación reticular de un automorfismo monomial \(h\) del código se desarrollan en la sección~\ref{paper:24-orientaciones}. Con todas las marcas transportadas,

\[
g_{hu}=\widetilde h\,g_u\,\widetilde h^{-1}.
\]

Este \(C_3\) es un grupo permisible realmente entrelazado con las fibras APP y la estructura excepcional. No se lo identifica aquí con \(C_9\) ni con su cociente por una mera relación de órdenes. La elevación al retículo tampoco se reemplaza por la matriz ternaria sin su acción sobre las fibras \(A_2\).

\subsection{Refinamiento de la energía y de su diferencial}
\label{paper:refinamiento-energia}

El operador de energía de \eqref{vii:eq:energia-memoria-discreta} tiene la forma \(D_{9,m}^*W_mD_{9,m}\), con \(U_a\) transportando orientación y acarreo. Los dominios de \cref{vii:subsec:funcional-memoria} son:

\[
\mathcal H_m=\bigoplus_vH_v,\quad
\mathcal K_m=\bigoplus_{a\in\mathcal E_m^+}H_{t(a)},\quad
(D_mf)_a=f_{t(a)}-U_af_{s(a)},\quad
E_m=\langle D_mf,\mathsf W_mD_mf\rangle.
\]

Las aplicaciones fijas \(J_m\) y \(K_m\) satisfacen, para la familia \(C^1\) considerada,

\[
D_{m+1}(t)J_m=K_mD_m(t),\qquad
K_m^*\mathsf W_{m+1}(t)K_m=\mathsf W_m(t).
\]

El \cref{vii:thm:refinamiento-respuesta} prueba \(E_{m+1}(J_mf)=E_m(f)\) por sustitución, y la igualdad de diferenciales se obtiene diferenciando ambas identidades. Para un refinamiento dependiente del parámetro se conserva también el término \(D_{m+1}\dot J_mf\).

Con \(v_a=U_af_s\) y \(q=\mathsf W_mD_mf\), la respuesta a \(\dot U_a=A_aU_a\) es

\[
\mathfrak j_m(A)=-2\operatorname{Re}\sum_a\langle q_a,A_av_a\rangle,
\quad \mathcal J_a=v_aq_a^*-q_av_a^*.
\]

El producto de transportes \(U_\gamma=U_N\cdots U_1\) cumple

\[
\dot U_\gamma U_\gamma^{-1}
=\sum_{k=1}^N B_kA_kB_k^{-1},\qquad B_k=U_N\cdots U_{k+1}.
\]

Esta fórmula retiene el orden de cada incidencia y transporta contragredientemente su covector. El corolario~\ref{vii:cor:refinamiento-corriente-conexion} compone ese diferencial con las variaciones de conexión y prueba, para los refinamientos compatibles de la acción,

\[
s_m=P_m^{\mathsf T}s_{m+1},\qquad
M_m\sigma_m=P_m^{\mathsf T}M_{m+1}\sigma_{m+1}.
\]

Aquí \(P_m\) lleva variaciones de conexión y \(M_m\) es el emparejamiento celular; no son el proyector de media nula ni la matriz de Hadamard.

\subsubsection{Ensamblaje con la incidencia}

En una realización de fibras \(W_H\) cuyas acciones de marco sean \(\widehat\rho_H(g_a)\), la conjugación por \(B\) lleva cada transporte, campo y peso a \(E_{\rm inc}\). Por el entrelazamiento demostrado en la sección~\ref{nuc04:kwitt}, los operadores de diferencias satisfacen \(D_{\rm inc}B_{\rm vertices}=B_{\rm edges}D_H\). Transportando también \(J_m\), \(K_m\) y \(\mathsf W_m\), los dos cuadrados de refinamiento anteriores siguen siendo los mismos cuadrados conjugados. Por sustitución conservan la energía; para \(B\) fijo, conservan además el diferencial y los covectores.

La ley de balance y sus covectores se transportan a la representación excepcional conservando su compatibilidad con el refinamiento. Cada aplicación utiliza los transportes de la dinámica especificada y los subespacios que preservan. El dominio de las variaciones está definido en \cref{vii:subsec:funcional-memoria}.
